% GENERATED FILE. Edit canonical lessons, then run npm run generate:books.
% canonical-source-hash: fnv1a64:b9d56fa998690d5e
% canonical-lessons: ML-C495-kotumuti, ML-C495-malancerivu, ML-C495-curam, ML-C495-pulmetu, ML-C495-tanal

\chapter{In the Hills}
\label{ch:ml-in-the-hills}
\hlchaptermodality{495}

\begin{chapteropening}
\textbf{By the end of this chapter:} \emph{I can say a peak, a hillside, a mountain pass, a meadow and shade.}

Complete the last lesson of chapter 495: I can say a peak, a hillside, a mountain pass, a meadow and shade.
\end{chapteropening}

\section[\texorpdfstring{\ml{കൊടുമുടി} (koṭumuṭi)}{കൊടുമുടി (koṭumuṭi)}]{\ml{കൊടുമുടി} (koṭumuṭi) — a peak; a summit}

\label{lesson:ML-C495-kotumuti}



\begin{quote}
Before the new one: say the Malayalam for to count and check, then the Malayalam for to sell off.
\end{quote}

\subsection*{You'll want to know: \ml{കൊടുമുടി}}
\textbf{\ml{കൊടുമുടി}} — \emph{koṭumuṭi} — \textquotedblleft{}a peak; a summit\textquotedblright{}.

\begin{quote}
\ml{മഞ്ഞുകാലത്ത്} \ml{ആ} \ml{കൊടുമുടിയിൽ} \ml{മഞ്ഞ്} \ml{വീഴും}.
\end{quote}

\begin{quote}
\emph{maññukālattŭ ā koṭumuṭiyil maññŭ vīḻuṁ.}
\end{quote}

\begin{quote}
In winter snow falls on that peak.
\end{quote}

Also the peak of anything: the top of fame or of a career.

\begin{cousinweb}
\textbf{\ml{കൊടുമുടി}} (\emph{koṭumuṭi}) is a Malayalam compound whose second part is \ml{മുടി}, the hair or crown of the head (DEDR 4921); the origin of the first part is not settled.
\end{cousinweb}

\begin{grammarlens}[title={What you've built}]
One more word for getting around.
\end{grammarlens}

\subsection*{Your turn}
\begin{itemize}
\raggedright
  \item \emph{Say it:} \emph{koṭumuṭi}
  \item \emph{Say it:} \emph{koṭumuṭi}, once more
  \item \emph{Say it:} \emph{maññukālattŭ ā koṭumuṭiyil maññŭ vīḻuṁ.} — In winter snow falls on that peak.
  \item \emph{Recall:} say the Malayalam for to count and check, then the Malayalam for to sell off, then say \emph{koṭumuṭi} again
\end{itemize}

\begin{tcolorbox}[breakable,colback=teal!4,colframe=teal!35!black,title={Before you move on}]
What does \ml{കൊടുമുടി} mean? (\textquotedblleft{}A peak; a summit\textquotedblright{}.) Say it once more.
\end{tcolorbox}
\section[\texorpdfstring{\ml{മലഞ്ചെരിവ്} (malañcerivŭ)}{മലഞ്ചെരിവ് (malañcerivŭ)}]{\ml{മലഞ്ചെരിവ്} (malañcerivŭ) — a hillside; a slope}

\label{lesson:ML-C495-malancerivu}



\begin{quote}
Before the new one: say the Malayalam for to sell off, then the Malayalam for a peak.
\end{quote}

\subsection*{You'll want to know: \ml{മലഞ്ചെരിവ്}}
\textbf{\ml{മലഞ്ചെരിവ്}} — \emph{malañcerivŭ} — \textquotedblleft{}a hillside; a slope\textquotedblright{}.

\begin{quote}
\ml{മലഞ്ചെരിവിൽ} \ml{ധാരാളം} \ml{തേയിലച്ചെടികൾ} \ml{ഉണ്ട്}.
\end{quote}

\begin{quote}
\emph{malañcerivil dhārāḷaṁ tēyilacceṭikaḷ uṇṭŭ.}
\end{quote}

\begin{quote}
There are many tea bushes on the hillside.
\end{quote}

\ml{മല}, a hill, plus \ml{ചെരിവ്}, a slope; \ml{ചരിവ്} and \ml{ചെരിവ്} are two spellings of one word.

\begin{cousinweb}
\textbf{\ml{മലഞ്ചെരിവ്}} (\emph{malañcerivŭ}) is a native Dravidian compound of \ml{മല}, a hill (DEDR 4742), and \ml{ചെരിവ്}, a slope (DEDR 2360).
\end{cousinweb}

\begin{grammarlens}[title={What you've built}]
One more word for getting around.
\end{grammarlens}

\subsection*{Your turn}
\begin{itemize}
\raggedright
  \item \emph{Say it:} \emph{malañcerivŭ}
  \item \emph{Say it:} \emph{malañcerivŭ}, once more
  \item \emph{Say it:} \emph{malañcerivil dhārāḷaṁ tēyilacceṭikaḷ uṇṭŭ.} — There are many tea bushes on the hillside.
  \item \emph{Recall:} say the Malayalam for to sell off, then the Malayalam for a peak, then say \emph{malañcerivŭ} again
\end{itemize}

\begin{tcolorbox}[breakable,colback=teal!4,colframe=teal!35!black,title={Before you move on}]
What does \ml{മലഞ്ചെരിവ്} mean? (\textquotedblleft{}A hillside; a slope\textquotedblright{}.) Say it once more.
\end{tcolorbox}
\section[\texorpdfstring{\ml{ചുരം} (curaṁ)}{ചുരം (curaṁ)}]{\ml{ചുരം} (curaṁ) — a mountain pass; a ghat road}

\label{lesson:ML-C495-curam}



\begin{quote}
Before the new one: say the Malayalam for a peak, then the Malayalam for a hillside.
\end{quote}

\subsection*{You'll want to know: \ml{ചുരം}}
\textbf{\ml{ചുരം}} — \emph{curaṁ} — \textquotedblleft{}a mountain pass; a ghat road\textquotedblright{}.

\begin{quote}
\ml{ബസ്} \ml{പതുക്കെ} \ml{ചുരം} \ml{കയറി}.
\end{quote}

\begin{quote}
\emph{basŭ patukke curaṁ kayaṟi.}
\end{quote}

\begin{quote}
The bus slowly climbed the ghat road.
\end{quote}

The road that bends up through the hills; its hairpin bends are \ml{ഹെയർപിൻ} \ml{വളവുകൾ}.

\begin{cousinweb}
\textbf{\ml{ചുരം}} (\emph{curaṁ}) is native Dravidian, \textquotedblleft{}a narrow and difficult path, a pass\textquotedblright{} (DEDR 2683).
\end{cousinweb}

\begin{grammarlens}[title={What you've built}]
One more word for getting around.
\end{grammarlens}

\subsection*{Your turn}
\begin{itemize}
\raggedright
  \item \emph{Say it:} \emph{curaṁ}
  \item \emph{Say it:} \emph{curaṁ}, once more
  \item \emph{Say it:} \emph{basŭ patukke curaṁ kayaṟi.} — The bus slowly climbed the ghat road.
  \item \emph{Recall:} say the Malayalam for a peak, then the Malayalam for a hillside, then say \emph{curaṁ} again
\end{itemize}

\begin{tcolorbox}[breakable,colback=teal!4,colframe=teal!35!black,title={Before you move on}]
What does \ml{ചുരം} mean? (\textquotedblleft{}A mountain pass; a ghat road\textquotedblright{}.) Say it once more.
\end{tcolorbox}
\section[\texorpdfstring{\ml{പുൽമേട്} (pulmēṭŭ)}{പുൽമേട് (pulmēṭŭ)}]{\ml{പുൽമേട്} (pulmēṭŭ) — a meadow; a grassland}

\label{lesson:ML-C495-pulmetu}



\begin{quote}
Before the new one: say the Malayalam for a hillside, then the Malayalam for a mountain pass.
\end{quote}

\subsection*{You'll want to know: \ml{പുൽമേട്}}
\textbf{\ml{പുൽമേട്}} — \emph{pulmēṭŭ} — \textquotedblleft{}a meadow; a grassland\textquotedblright{}.

\begin{quote}
\ml{മലമുകളിലെ} \ml{പുൽമേട്ടിൽ} \ml{പശുക്കൾ} \ml{മേഞ്ഞു}.
\end{quote}

\begin{quote}
\emph{malamukaḷile pulmēṭṭil paśukkaḷ mēññu.}
\end{quote}

\begin{quote}
Cows grazed on the meadow on the hilltop.
\end{quote}

\ml{പുല്ല്}, grass, plus \ml{മേട്}, high ground.

\begin{cousinweb}
\textbf{\ml{പുൽമേട്}} (\emph{pulmēṭŭ}) is a native Dravidian compound of \ml{പുല്ല്}, grass (DEDR 4300), and \ml{മേട്}, a rise or mound of land (DEDR 5058).
\end{cousinweb}

\begin{grammarlens}[title={What you've built}]
One more word for getting around.
\end{grammarlens}

\subsection*{Your turn}
\begin{itemize}
\raggedright
  \item \emph{Say it:} \emph{pulmēṭŭ}
  \item \emph{Say it:} \emph{pulmēṭŭ}, once more
  \item \emph{Say it:} \emph{malamukaḷile pulmēṭṭil paśukkaḷ mēññu.} — Cows grazed on the meadow on the hilltop.
  \item \emph{Recall:} say the Malayalam for a hillside, then the Malayalam for a mountain pass, then say \emph{pulmēṭŭ} again
\end{itemize}

\begin{tcolorbox}[breakable,colback=teal!4,colframe=teal!35!black,title={Before you move on}]
What does \ml{പുൽമേട്} mean? (\textquotedblleft{}A meadow; a grassland\textquotedblright{}.) Say it once more.
\end{tcolorbox}
\section[\texorpdfstring{\ml{തണൽ} (taṇal)}{തണൽ (taṇal)}]{\ml{തണൽ} (taṇal) — shade}

\label{lesson:ML-C495-tanal}



\begin{quote}
Before the new one: say the Malayalam for a mountain pass, then the Malayalam for a meadow.
\end{quote}

\subsection*{You'll want to know: \ml{തണൽ}}
\textbf{\ml{തണൽ}} — \emph{taṇal} — \textquotedblleft{}shade\textquotedblright{}.

\begin{quote}
\ml{ഉച്ചയ്ക്ക്} \ml{ഞങ്ങൾ} \ml{ഒരു} \ml{വലിയ} \ml{മരത്തിന്റെ} \ml{തണലിൽ} \ml{ഇരുന്നു}.
\end{quote}

\begin{quote}
\emph{uccaykkŭ ñaṅṅaḷ oru valiya marattinṟe taṇalil irunnu.}
\end{quote}

\begin{quote}
At noon we sat in the shade of a big tree.
\end{quote}

Shade is a kind of coolness: the word comes from the same root as \ml{തണുപ്പ്}.

\begin{cousinweb}
\textbf{\ml{തണൽ}} (\emph{taṇal}) is native Dravidian, from the root of \ml{തണുപ്പ്}, cold; DEDR lists it as shade, a shady spot (DEDR 3045).
\end{cousinweb}

\begin{grammarlens}[title={What you've built}]
One more word for getting around.
\end{grammarlens}

\subsection*{Your turn}
\begin{itemize}
\raggedright
  \item \emph{Say it:} \emph{taṇal}
  \item \emph{Say it:} \emph{taṇal}, once more
  \item \emph{Say it:} \emph{uccaykkŭ ñaṅṅaḷ oru valiya marattinṟe taṇalil irunnu.} — At noon we sat in the shade of a big tree.
  \item \emph{Recall:} say the Malayalam for a mountain pass, then the Malayalam for a meadow, then say \emph{taṇal} again
\end{itemize}

\begin{tcolorbox}[breakable,colback=teal!4,colframe=teal!35!black,title={Before you move on}]
What does \ml{തണൽ} mean? (\textquotedblleft{}Shade\textquotedblright{}.) Say it once more.
\end{tcolorbox}
